\chapter{Einleitung}

Die Messung von Durchflüssen ist ein fundamentaler Aspekt in zahlreichen Bereichen der Technik. Die präzise Erfassung von Stoffströmen spielt eine entscheidende Rolle für Effizienz, Sicherheit und Qualitätskontrolle, denn die Regelung und Bilanzierung der Stoffströme kann nur mit genauen Werten durchgeführt werden. In diesem Praktikumstermin wurden zwei Versuche, mit dem Ziel, die Durchflüsse verschiedener Medien kennenzulernen, durchgeführt. Der erste mit dem inkompressiblen Medium Wasser, der zweite mit dem kompressiblen Medium Luft. Die Vorgaben zur Versuchsdurchführung und Hintergründe dafür sind durch die Praktikumsanleitung gegeben. \cite{waermeatlas} \cite{CanJour}